\chapter{結論與未來工作}
本篇論文基於\cite{9026293}建立的腳踏車騎乘者生理資訊、環境資訊收集，再加上電子變速系統，對其系統進行關於採樣頻率與系統穩定度方面的調整，並且利用調整後的平台對校內5條不同騎乘路線、與校外路線進行訓練資料的收集，另外收集了一筆校內、與校外資料作為我們的測試路線。使用數個機器學習、深度學習的方法對該資料集進行自車檔位的預測，使用的方法包含了神經網路、RandomForests、DeepForest、XGBoost、LightGBM、CATBoost這幾個，並且對這些方法進行調整優化，對所有的結果進行分析與討論，最後我們認為神經網路方法依舊無法有效的對表格類資料實行有效預測，其餘的方法中RandomForest、XGBoost、LightGBM、CATBoost得益於可以反覆進行學習因而得到了最令人滿意的效果，其中又以XGBoost與CATBoost的預測最為平順且符合測試資料的騎乘習慣。另外我們建立一個網頁平台，用於資料的儲存、查詢，與簡易的分析比較。供使用者了解自己的騎乘狀況，並確認該系統的預測與真實的狀況是否吻合。


本論文的主要貢獻如下:
\begin{enumerate}
	\item 調整先前建立平台的資料收集頻率，從1分鐘5-6筆(約10-15秒一筆)，提升至1秒3筆資料。
	\item 收集建立了一個包含踩踏壓力、IMU、心率、車速、時間、GPS 座標資訊、環境溫度、圖資，當下檔位狀態資料集。該資料集包含了約10,000筆訓練資料，各2,000筆校內、校外測試資料。
	\item 使用多種方法對該數據集進行分析，並探討這些方法對於我們資料集的實用性。
	\item 建立一種兩階段的訓練方法，對時序相關資料進行參考。
	\item 建立一個可供使用者管理並查詢騎乘狀況的網頁平台。
\end{enumerate}

未來可改善的方向，首先從我們的實驗中看到校外資料尚未達到飽和還有一些提升的空間，因此可以先著手於校外資料集的增加，資料增加方式可以使用收集資料的方式，或是最近有提出的GAN方法對我們的資料集進實施增強。另外目前實驗都是在個人電腦上(PC)運行的，可將預測階段移植到樹梅派上實行真正的即時預測與換檔。最後希望可以透過更換平台或是其他方式改善傳輸延遲的問題即時將資料上傳。另外我們有設計了一個分數評估的標準(score)，因此未來可以考慮使用該分數作為訓練時的參考依據，現在沒使用該方法主要是因為我們的分數目前僅適合，作為驗證的分數評估，並不適合作為目標函數使用，所以未來可以先將分數評估的方式調整成適合做為目標函數的形式後，已進行利用。